% =====================================================================
%  CL4FMAgents 2026 submission -- BODY v2 (after FC review pass)
%  Replaces everything from \documentclass down to \bibliography.
%  KEEP your existing \appendix ... \end{document} unchanged below.
% =====================================================================
\documentclass{article}

% Double-blind workshop (anonymous). If CL4FMAgents turns out to be
% SINGLE-blind, swap the two lines below and put the real author block back.
\usepackage[dblblindworkshop]{neurips_2026}
% \usepackage[sglblindworkshop]{neurips_2026}

\usepackage[utf8]{inputenc}
\usepackage[T1]{fontenc}
\usepackage[hidelinks]{hyperref}   % hidelinks = no coloured boxes around links
\usepackage{url}
\usepackage{booktabs}
\usepackage{amsfonts}
\usepackage{amsmath}
\usepackage{amssymb}
\usepackage{nicefrac}
\usepackage{microtype}
\usepackage{xcolor}
\usepackage{graphicx}
\usepackage{subcaption}

\title{Cross-Modal Forgetting and Transfer in Parameter-Efficient\\
Post-Training of a Unified Multimodal Model}

\workshoptitle{NeurIPS 2026 Workshop on Continual Learning in the Era of
Foundation Models and Embodied Agents}

\author{%
  Anonymous Authors \\
}

\begin{document}
\maketitle

\begin{abstract}
Unified multimodal models serve image understanding and image generation from a
single shared backbone, which makes post-training them a continual-learning
problem with an unusual shape: acquiring a new capability can degrade
capabilities that live in a \emph{different modality} of the same weights. We
measure this cross-modal catastrophic forgetting for parameter-efficient
post-training (LoRA) of Janus-Pro-7B. We adapt each direction in turn on five
targets chosen by how poorly the base model already handles them, two visual
styles (anime, line sketches) and three text targets (poetry, medical abstracts,
Welsh), and after each adaptation we evaluate every capability we did \emph{not}
train: MMLU, POPE, MME and WikiText perplexity, free-form generation quality
(MAUVE for text, a distributional distance for images), and the weight-space size
of the update. Across LoRA ranks $2$ to $128$ and KL weights
$\lambda\in[0,1.5]$, the untrained capabilities stay at base level within
evaluation noise, even under a strong shift toward a language the base barely
modelled (held-out Welsh perplexity $-64\%$). Where any cost appears at all it is
small and confined to grounded visual understanding, and only at the strongest,
least-regularized settings; we are explicit that single-seed evidence cannot
establish a null result, and we say what would. The same locality has a cost we
quantify: the acquired language does \emph{not} transfer to the visual pathway,
leaving Welsh visual question answering at chance. The KL weight controls how far
the model drifts, but it does not control whether what is learned reaches the
other modality.
\end{abstract}

\section{Introduction}
Continual learning asks how a model can acquire a new capability without
destroying the ones it already has \citep{mccloskey1989catastrophic,
french1999catastrophic}. Foundation models have made the question practical
rather than academic: post-training is now the normal way to specialize a
released checkpoint, and every such update risks eroding the pretrained
competence that motivated using the checkpoint in the first place
\citep{kirkpatrick2017ewc, li2017lwf}.

Unified multimodal models give this question an unusual shape. Models such as
Chameleon \citep{chameleon2024}, the 4M family \citep{bachmann20244m21},
Unified-IO~2 \citep{lu2024unifiedio2} and Janus \citep{wu2024janus, januspro}
serve both understanding (image$+$text to text) and generation (text to image)
from a single set of weights. When post-training specializes one capability, the
capabilities at risk are not only the neighbouring tasks of the same modality:
they include tasks in a \emph{different} modality that happen to share the same
parameters. We call this cross-modal forgetting. It cannot arise in a task
specialist, such as a text-to-image diffusion model, because there is no other
modality to preserve; it is the natural stress test of the ``one backbone, many
capabilities'' premise.

Parameter-efficient fine-tuning, LoRA in particular \citep{hu2022lora}, is the
default post-training tool and is often assumed to be safe by construction
because it moves few parameters. The assumption is under-tested in exactly the
setting where it matters most. Prior work bounds the problem without answering
it: LoRA forgets less than full fine-tuning while also learning less
\citep{biderman2024lora}, and increasing the LoRA rank beyond a small value buys
little additional task performance \citep{zhang2024scalingmeetsfinetuning}, but
both are measured within a single modality. Work that does address the
understanding/generation tension in unified models resolves it by changing the
model, routing among several adapters or adding a mixture-of-experts layer
\citep{dou2024loramoe}; such designs presuppose that plain LoRA interferes, but
none of them measures how much it actually does. That measurement is what we
provide.

Measuring it credibly requires care about the target. If the fine-tuning target
is something the base model already handles well, the adapter barely has to move
the weights, and finding ``no forgetting'' says nothing: a weak intervention
cannot damage anything, so the result is guaranteed in advance and would be
reported all the same. We therefore select every target by how poorly the base
already handles it, on a scale that is comparable across targets, and we report
how large the resulting weight update is, so that a reader can check that the
intervention was real before believing that nothing broke.

We report two findings on Janus-Pro-7B. First, post-training is strongly
localized: over five targets, LoRA ranks $2$ to $128$ and KL weights
$\lambda\in[0,1.5]$, the capabilities we did not train stay at base level within
evaluation noise, on the four benchmarks, on free-form generation, and in the
other modality's generative distribution. Second, the same localization has a
cost: a language acquired in the text stream is not usable through the visual
pathway at all. Isolation is what protects the untrained capabilities, and it is
also what prevents the new one from reaching them.

\paragraph{Contributions.}
\textbf{(i)} We show that claims of the form ``PEFT does not forget'' are only
meaningful when the adaptation target is demonstrably outside the base model's
existing competence, and we make that condition explicit and measurable: every
target is scored by the base model's perplexity on it relative to a fixed
control, and every adaptation is accompanied by the size of the weight update it
induced.
\textbf{(ii)} Under that condition, the untrained capabilities are preserved
across the whole range of adaptation strength we can reach, including under a
strong distributional shift toward a language the base barely modelled (held-out
perplexity $40.7\to14.6$), and beyond likelihood scoring: free-form English
generation is preserved (MAUVE) and the image-generation distribution is
essentially unmoved (Gen-CMMD $\sim\!1.7\times10^{-3}$). We locate where a cost
first appears, in grounded visual understanding at the strongest settings, and
report it with its uncertainty rather than as a demonstrated failure.
\textbf{(iii)} A negative transfer result that we argue is the flip side of the
same mechanism: the acquired language improves text-to-image faithfulness only
marginally ($+0.47\pm0.25$ CLIPScore, $3\%$ of the English--Welsh gap) and leaves
Welsh visual question answering at chance, where the model turns out not to be
reading the Welsh question at all.
\textbf{(iv)} Each adapter changes $0.6$ to $1.9\%$ of the backbone's weights,
concentrated in the MLP and sparing attention. The target that needs the smallest
update is not the one furthest from the base: Welsh reaches the largest
perplexity reduction of the three text targets while moving the weights least
($0.6\%$), suggesting the capability is unlocked rather than built.

\section{Related Work}
\paragraph{Continual learning and anti-forgetting objectives.}
Continual learning is classically framed as a tension between retaining what a
model can already do and acquiring something new, the stability--plasticity
dilemma \citep{mermillod2013stability}, which remains the organizing question of
the continual-learning literature on deep networks \citep{wang2024clsurvey}.
Classical remedies anchor the model to a reference, in weight space, L2-SP
\citep{xuhong2018l2sp} and its Fisher-weighted form EWC \citep{kirkpatrick2017ewc},
or in function space, distillation \citep{hinton2015distilling} and
Learning-without-Forgetting \citep{li2017lwf}. We use a function-space reference
and, following standard continual-learning practice \citep{lopezpaz2017gem},
report not only what is retained but also whether what is learned transfers to a
capability the model was not trained on.

\paragraph{PEFT and forgetting.}
Adapters \citep{houlsby2019parameter} and LoRA \citep{hu2022lora} update a small
fraction of parameters. LoRA's adaptation/forgetting trade-off is comparatively
benign: it learns less and forgets less than full fine-tuning
\citep{biderman2024lora}, and raising the rank beyond a small value buys little
\citep{zhang2024scalingmeetsfinetuning}. A complementary line mitigates
forgetting by routing among multiple LoRAs \citep{dou2024loramoe, wu2024mole,
li2024mixlora}. These works study text-only or late-fusion models and treat
forgetting within a single modality.

\paragraph{Unified multimodal models.}
Fully-tokenized native models \citep{chameleon2024, bachmann20244m21,
lu2024unifiedio2} process all modalities in one sequence. Janus and Janus-Pro
\citep{wu2024janus, januspro} decouple the visual encoders for understanding and
generation while keeping a single shared backbone. PEFT has been applied to such
models \citep{chern2024anole, serra2024narrowgate, wang2025vora}, but as a means
to an end rather than to characterize cross-capability interference. We neither
modify the architecture nor propose a new adapter: we measure what plain LoRA
does to the capabilities it was not trained on.

\section{Setup}
\paragraph{Model, directions, and adapter.}
We use Janus-Pro-7B \citep{januspro}: a SigLIP encoder for understanding and a
VQ-VAE (codebook $16{,}384$, $576$ tokens per image) for generation, over one
shared LLaMA-style transformer. It exposes three capability directions:
text-to-text language modelling, image$+$text understanding, and text-to-image
generation. LoRA is injected \emph{only} into the shared backbone
($q,k,v,o,\texttt{gate},\texttt{up},\texttt{down}$), with $\alpha=2r$ fixed
across the rank sweep; the encoders and heads are frozen. The generation forward
traverses only the language-model tensors, while understanding additionally
engages SigLIP.

\paragraph{Controlling how far the adapter drifts: a KL term to the frozen base.}
For both directions we optionally hold the adapter close to the frozen base.
Writing $p_{\text{adapter}}(\cdot\mid x_{<t})$ for the adapted next-token
distribution and $p_{\text{base}}(\cdot\mid x_{<t})$ for the \emph{same} model
with the adapter disabled, the training loss is
\begin{equation}
\mathcal{L}=\underbrace{\mathrm{CE}\big(p_{\text{adapter}},y\big)}_{\text{acquire the target}}
+\lambda\,\frac{1}{T}\sum_{t=1}^{T}\underbrace{\mathrm{KL}\!\big(p_{\text{adapter}}(\cdot\mid x_{<t})\,\big\|\,\mathrm{sg}[p_{\text{base}}(\cdot\mid x_{<t})]\big)}_{\text{stay close to the base}} .
\label{eq:kl}
\end{equation}
The reference is obtained for free by disabling the adapter, so no separate
teacher is stored, and it is held constant by a stop-gradient. The single KL
weight $\lambda$ sets how much the adapter may drift: $\lambda=0$ is plain
fine-tuning, larger $\lambda$ keeps the adapter near the base. We use a
function-space reference because the quantity to preserve is the model's output
\emph{distribution}; an ablation against L2-SP, logit-MSE and the reconstruction
objective of Anole \citep{chern2024anole} at matched update size supports the
choice (Appendix~\ref{app:reg}).

\paragraph{What we evaluate after adaptation.}
We evaluate the directions we did not train with four standard benchmarks scored
by log-likelihood only, hence deterministic and judge-free: WikiText-2 perplexity
\citep{wikitext}, MMLU \citep{mmlu}, POPE \citep{pope} and MME \citep{mme}. The
base POPE accuracy ($0.877$) matches the published Janus-Pro-7B value ($0.874$),
validating the harness. Because likelihood benchmarks never exercise free-form
decoding, we add two distributional generation metrics: MAUVE
\citep{pillutla2021mauve} for text, and, for the image direction under a text
adapter, the CMMD \citep{cmmd} between the adapter's and the base's generations
on a fixed prompt set (Gen-CMMD).

\paragraph{What we mean by retention.}
We call a capability \emph{retained} when the adapter's score on it is within
evaluation noise of the base's. We apply this symmetrically: a large movement is
a behavioural change whatever its sign, so the cases where a score moves
substantially in the \emph{favourable} direction (general-text perplexity, and
MAUVE-vs-base for the medicine adapter) are flagged as changes and analysed
separately in \S\ref{sec:dw}, not counted as evidence of stability. Our
benchmarks also carry a scope limit we state once here: they were never selected
for sensitivity to cross-modal damage, because no prior protocol establishes what
such damage looks like on this model. The one condition in which we do observe a
score move far outside noise is a control (joint training on a small text corpus,
WikiText perplexity $13.71\!\to\!55.10$, Appendix~\ref{app:ablation}), which
establishes that at least the perplexity measure is not inert.

\paragraph{Choosing targets the base cannot already do.}
As argued in the introduction, certifying that nothing was forgotten is
uninformative if the target was already within the base model's reach. We
therefore score each candidate target by $\times$control: the base model's
perplexity on that corpus divided by its perplexity on a fixed English control.
On this common scale the three text targets are poetry (Project Gutenberg) at
$4.0\times$, Welsh (\texttt{wikimedia/wikipedia}, Welsh subset) at $2.3\times$,
and medical abstracts (\texttt{ccdv/pubmed-summarization}) at $1.2\times$; Welsh
was selected from a wider set of low-resource and constructed languages
(Appendix~\ref{app:langprobe}). The
three therefore differ in \emph{kind}, a register, a language and a domain,
rather than lying on one ordered axis of difficulty, and we will see
(\S\ref{sec:text}) that $\times$control predicts neither how much is acquired nor
how far the weights move. Poetry and medicine were scored in one probe run and
Welsh in another, with the same construction and control corpus. For images we
use a narrow anime style (\texttt{lambdalabs/naruto-blip-captions}, $1{,}141$
pairs) that the base renders weakly, and black-and-white line sketches
(\texttt{zoheb/sketch-scene}). Image acquisition is scored against real style
images with CMMD \citep{cmmd} as the headline, with KID \citep{kid}, DINOv2
\citep{dinov2}, CLIP-IQA \citep{clipiqa} and CLIP directional similarity
\citep{stylegannada} as complements. Full hyperparameters are in
Appendix~\ref{app:impl}. All numbers are single-seed; standard errors quantify
evaluation-set sampling noise at fixed weights (Appendix~\ref{app:impl}).

\section{Results}
We adapt one direction at a time and evaluate every other direction.
Section~\ref{sec:gen} adapts image generation and evaluates text, knowledge and
understanding; Section~\ref{sec:text} adapts language modelling and evaluates
knowledge, understanding and the image-generation distribution. Together the two
form the off-diagonal of a retention matrix over the model's directions.

\begin{figure}[t]
\centering
\begin{subfigure}{0.24\textwidth}
\includegraphics[width=\linewidth]{figs/janus_comparison/no_adapter/human_05_natural.png}
\caption*{\footnotesize Base}
\end{subfigure}
\begin{subfigure}{0.24\textwidth}
\includegraphics[width=\linewidth]{figs/janus_comparison/adapter_solo_gen/human_05_natural.png}
\caption*{\footnotesize Image LoRA}
\end{subfigure}
\begin{subfigure}{0.24\textwidth}
\includegraphics[width=\linewidth]{figs/janus_comparison/no_adapter/human_07_natural.png}
\caption*{\footnotesize Base}
\end{subfigure}
\begin{subfigure}{0.24\textwidth}
\includegraphics[width=\linewidth]{figs/janus_comparison/adapter_solo_gen/human_07_natural.png}
\caption*{\footnotesize Image LoRA}
\end{subfigure}
\caption{The intervention is real: style acquisition on two held-out prompts
(``a young man with blonde spiky hair'', ``an old wizard with a long white
beard''), base versus adapter at identical seeds. Full grids in
Appendix~\ref{app:grids}.}
\label{fig:teaser}
\end{figure}

\subsection{Adapting the generation direction}
\label{sec:gen}
We fine-tune generation to a final configuration ($r=48$, attention and MLP
LoRA, $\lambda=0.5$; the incremental ablation reaching it is in
Appendix~\ref{app:ablation}). The adapter acquires the style
(Figure~\ref{fig:teaser}).

Table~\ref{tab:rank} sweeps the rank from $2$ to $64$, i.e.\ the adapter's
capacity to change the shared backbone. Acquisition improves and saturates
(CMMD $0.139\to0.097$, directional similarity $0.143\to0.158$, knee around rank
$16$ to $32$), while every untrained capability stays at the base level across
the whole range: WikiText perplexity moves within $0.07$ of the base, MMLU, POPE
and MME within sampling noise. More capacity refines acquisition at no measurable
cost to the other modalities. We replicate this on a second, unrelated visual
style (line sketches, Appendix~\ref{app:sketch}); we flag now, rather than leave
to the appendix, that this second style produces the largest footprint of all
five targets and is one of the three settings analysed in \S\ref{sec:boundary}.

\begin{table}[t]
\caption{\textbf{Adapting image generation: rank sweep.} Acquisition (bold =
best) improves and saturates, while every untrained capability stays at the base
level at all ranks. Base POPE ($0.877$) matches the published value ($0.874$).
The base has no directional similarity by construction (DirSim measures the
change away from the base), and its CMMD to the style reference was not computed
in this sweep; Table~\ref{tab:sketch} reports a base CMMD for the second style.}
\label{tab:rank}
\centering\small
\begin{tabular}{lccccccc}
\toprule
& Base & $r{=}2$ & $r{=}4$ & $r{=}8$ & $r{=}16$ & $r{=}32$ & $r{=}64$ \\
\midrule
\multicolumn{8}{l}{\emph{Acquisition} (vs.\ real style images)}\\
CMMD $\downarrow$        & n/a & .139 & .124 & .126 & .105 & .108 & \textbf{.097} \\
CLIP DirSim $\uparrow$   & n/a & .143 & .142 & .143 & .154 & .153 & \textbf{.158} \\
\midrule
\multicolumn{8}{l}{\emph{Untrained capabilities}}\\
WikiText PPL $\downarrow$ & 13.71 & 13.66 & 13.70 & 13.76 & 13.74 & 13.78 & 13.76 \\
MMLU $\uparrow$           & .492 & .491 & .491 & .491 & .501 & .497 & .489 \\
POPE acc $\uparrow$       & .877 & .879 & .881 & .878 & .878 & .879 & .878 \\
MME $\uparrow$            & 1785 & 1777 & 1777 & 1786 & 1786 & 1796 & 1781 \\
\bottomrule
\end{tabular}
\end{table}

\subsection{Adapting the language-modelling direction}
\label{sec:text}
We now specialize text on three targets that differ in kind: a register
(poetry), a domain (medicine) and a language (Welsh). For each we report
acquisition together with the untrained capabilities and the image-generation
shift.

\paragraph{Sweeping the KL weight.}
Table~\ref{tab:textkl} sweeps $\lambda$ at rank $16$ on poetry. Acquisition is
strongest with little regularization and weakens monotonically as the adapter is
held closer to the base (poetry perplexity $24.5$ at $\lambda=0$ rising to
$26.7$ at $\lambda=1.5$), the expected trade between acquiring and preserving.
Knowledge and POPE stay flat. Crucially, the bottom row shows the \emph{other
modality}: Gen-CMMD stays flat and near zero ($\sim\!1.7\times10^{-3}$) at every
$\lambda$, including $\lambda=0$, where the text objective is entirely
unconstrained. Even with no KL term at all, specializing the text stream does not
measurably move the image-generation distribution. Appendix~\ref{app:t2itext}
shows the same conclusion qualitatively, and Appendix~\ref{app:textsweep} reports
the complementary rank sweep on this target.

\begin{table}[t]
\caption{\textbf{Adapting text (poetry): KL sweep at rank 16.} Acquisition (bold
= best) is strongest at low $\lambda$ and weakens as the KL term tightens. MMLU
and POPE stay flat; MME shows a small dip at $\lambda=0$ (see
\S\ref{sec:boundary}). \emph{Gen-CMMD} is the CMMD between the adapter's and the
base's image generations on a fixed prompt set ($N{=}128$): it stays near zero at
every $\lambda$, so the text adapter does not measurably shift image generation.
Text runs use a subsampled evaluation configuration, so absolute values are
comparable within the table, against the base column, not against
Table~\ref{tab:rank}.}
\label{tab:textkl}
\centering\small
\begin{tabular}{lcccccc}
\toprule
& Base & $\lambda{=}0$ & $\lambda{=}0.25$ & $\lambda{=}0.5$ & $\lambda{=}1$ & $\lambda{=}1.5$ \\
\midrule
\multicolumn{7}{l}{\emph{Acquisition} (target style)}\\
Poetry PPL $\downarrow$ & 32.7 & \textbf{24.5} & \textbf{24.5} & 24.9 & 25.9 & 26.7 \\
\midrule
\multicolumn{7}{l}{\emph{Untrained capabilities}}\\
WikiText PPL $\downarrow$ & 13.71 & 12.61 & \underline{12.50} & 12.51 & 12.75 & 12.93 \\
MMLU $\uparrow$           & .481 & .474 & .481 & .479 & .475 & .483 \\
POPE acc $\uparrow$       & .875 & .868 & .870 & .871 & .878 & .873 \\
MME $\uparrow$            & 1130 & 1101 & 1116 & 1128 & 1127 & 1128 \\
\midrule
\multicolumn{7}{l}{\emph{Shift in the other modality}}\\
Gen-CMMD $\downarrow$ ($\times10^{-3}$) & 0 & 1.73 & 1.71 & 1.73 & 1.77 & 1.65 \\
\bottomrule
\end{tabular}
\end{table}

\paragraph{A domain and a language.}
Poetry and medicine both lie inside the base's competence (poetic English,
medical English), so the absence of forgetting could still be attributed to a
weak intervention. We therefore add a target the base barely speaks: Welsh.
Table~\ref{tab:targets} reports both. Medicine is acquired strongly (held-out
PubMed perplexity $12.8\to8.4$; on the $\times$control scale the domain moves
from $1.2\times$, harder than generic English, to $0.79\times$, easier), and
Welsh more strongly still (held-out perplexity $40.7\to14.6$, $-64\%$); on Welsh
seeds the adapter sustains Welsh where the base slips into English or its
assistant persona (Appendix~\ref{app:welsh}). Qualitative examples for medicine
are in Appendix~\ref{app:medicine}. In both cases the untrained capabilities stay
at base level and Gen-CMMD stays at the same $\sim\!1.7\times10^{-3}$ order as
poetry.

Two observations sharpen the picture. First, \emph{distance from the base
predicts neither outcome}. On the common $\times$control scale Welsh ($2.3\times$)
sits between poetry ($4.0\times$) and medicine ($1.2\times$), yet it produces the
largest perplexity reduction of the three and needs the \emph{smallest} weight
change ($\lVert\Delta W\rVert/\lVert W\rVert=0.6\%$, versus $1.3\%$ for poetry
and $1.7\%$ for medicine), despite the weakest KL term ($\lambda=0.1$). How
unfamiliar a corpus looks to the base model therefore tells us little about how
much of the network has to move to fit it; the new language appears to be
unlocked rather than built, i.e.\ partly latent in the pretrained backbone.
Second, under the adapter, English prompts are still continued in English, so the
language is added without overwriting the existing one.

We are deliberately careful about how much this licenses. A $64\%$ perplexity
reduction and a MAUVE of $0.67$ against real Welsh establish that the output
distribution moves substantially toward the target language; they do not
establish competence. Our own samples (Appendix~\ref{app:welsh}) show the adapter
sustaining Welsh morphology and register where the base does not, while remaining
locally repetitive and often semantically empty under greedy decoding. We
therefore claim a strong distributional shift toward a previously unmodelled
language, not the acquisition of a usable one, and the retention result should be
read against that weaker premise.

\begin{table}[t]
\caption{\textbf{A distant domain and a new language.} Medicine and Welsh are
each acquired strongly while every untrained capability stays at base level and
image generation is essentially unchanged. The two targets use different
evaluation configurations (full for medicine, subsampled for Welsh), so their
base columns differ in absolute value; each comparison is within its own pair of
columns. Errors are one standard error from evaluation-set resampling. Entries
marked n/a are metrics defined only for the other target.}
\label{tab:targets}
\centering\small
\setlength{\tabcolsep}{5pt}
\begin{tabular}{lcccc}
\toprule
& \multicolumn{2}{c}{Medicine ($r{=}32$, $\lambda{=}0.25$)}
& \multicolumn{2}{c}{Welsh ($r{=}32$, $\lambda{=}0.1$)} \\
\cmidrule(lr){2-3}\cmidrule(lr){4-5}
& Base & LoRA & Base & LoRA \\
\midrule
\multicolumn{5}{l}{\emph{Acquisition}}\\
PubMed PPL $\times$control $\downarrow$ & 1.23$\times$ & \textbf{0.79$\times$} & n/a & n/a \\
Welsh held-out PPL $\downarrow$         & n/a & n/a & 40.7 & \textbf{14.6} \\
$\lVert\Delta W\rVert/\lVert W\rVert$   & --- & 1.7\% & --- & 0.6\% \\
\midrule
\multicolumn{5}{l}{\emph{Untrained capabilities}}\\
WikiText PPL $\downarrow$ & 13.71\,$\pm$.21 & 13.97\,$\pm$.21 & 13.71\,$\pm$.21 & 12.07\,$\pm$.18 \\
MMLU $\uparrow$           & .492\,$\pm$.013 & .495\,$\pm$.013 & .481\,$\pm$.021 & .479\,$\pm$.016 \\
POPE acc $\uparrow$       & .877\,$\pm$.006 & .877\,$\pm$.006 & .875\,$\pm$.010 & .876\,$\pm$.010 \\
MME $\uparrow$            & 1785\,$\pm$31   & 1766\,$\pm$32   & 1130\,$\pm$42   & 1124\,$\pm$42 \\
\midrule
\multicolumn{5}{l}{\emph{Shift in the other modality}}\\
Gen-CMMD $\downarrow$ ($\times10^{-3}$) & 0 & 1.77 & 0 & 1.73 \\
\bottomrule
\end{tabular}
\end{table}

\subsection{Beyond likelihood: free-form generation}
The four benchmarks score multiple-choice and log-likelihood items, never
free-form decoding, so a generative degradation could hide behind flat
likelihood numbers. We close that gap with MAUVE \citep{pillutla2021mauve}, the
text counterpart of our image CMMD (Table~\ref{tab:mauve}). Three readings
follow. \emph{Cross-modal}: the image adapters generate free English essentially
identical to the base (MAUVE-vs-base $0.95$/$0.92$), the generative analogue of
their near-zero Gen-CMMD. \emph{Text, with one caveat}: no adapter's free English
falls below the base against human text, and the text adapters raise it
($0.60$--$0.77$), so none of them degrades English relative to the human
reference. We stress what this does \emph{not} say. A MAUVE-vs-base of $0.21$
(medicine) is a large distributional change: that adapter's free English is far
from the base's, even while being closer to human text. Under our definition
(\S3) this is a behavioural change, not retention. What the two columns jointly
rule out is \emph{degradation}, which is the weaker and the honest claim.
\emph{Acquisition}: the text adapters move far closer to the target distribution
than the base (poetry $0.45\to0.80$, medical abstracts $0.31\to0.91$, Welsh
$0.14\to0.67$), confirming acquisition at the generation-distribution level and
not only in perplexity. Appendix~\ref{app:textqual} shows the corresponding
poetry continuations.

\begin{table}[t]
\caption{\textbf{Free-form generation (MAUVE)} \citep{pillutla2021mauve}
($[0,1]$, higher = closer; $N\!\le\!2000$, GPT-2-large features). \emph{English
vs.\ base}: image adapters are indistinguishable from the base. \emph{English
vs.\ human}: no adapter falls below the base ($0.38$); text adapters raise it.
\emph{Target acquisition}: text adapters move far closer to the target. Base
continuations are coherent English, not degenerate (Appendix~\ref{app:mauve}).}
\label{tab:mauve}
\centering\small
\begin{tabular}{lccc}
\toprule
Adapter & English & English & Target acq. \\
        & vs.\ human & vs.\ base & base\,$\to$\,adapter \\
\midrule
Base (reference)        & 0.38 & --- & --- \\
Anime, image (Final)    & 0.52 & 0.95 & --- \\
Sketch, image           & 0.46 & 0.92 & --- \\
Poetry                  & 0.70 & 0.83 & 0.45\,$\to$\,0.80 \\
Medicine                & 0.60 & 0.21 & 0.31\,$\to$\,0.91 \\
Welsh                   & 0.77 & 0.68 & 0.14\,$\to$\,0.67 \\
\bottomrule
\end{tabular}
\end{table}

\subsection{Where a cost first appears}
\label{sec:boundary}
Reporting only preserved capabilities would say nothing about where the limit
is, so we deliberately searched for the settings that break it. Three
independent settings, each at the strongest end of adaptation, show a movement
in the same direction and in the same measure, grounded visual understanding
(MME), while knowledge and language stay flat:
(i) at $\lambda=0$, the unregularized poetry adapter, MME $1130\to1101$
(Table~\ref{tab:textkl}), which the KL term removes;
(ii) pushing poetry to rank $128$, well past the acquisition knee, buys almost
nothing ($23.8\to23.2$) while MME moves $1122\to1105$ and POPE $0.877\to0.867$;
(iii) the hardest visual style, line sketches, is the only target for which the
footprint is uniformly larger than the others, most visibly on MME
($1785\to1713$; Appendix~\ref{app:sketch}).

We state plainly what this is and is not. Each movement individually lies within
roughly one to two standard errors of the difference, and all runs are
single-seed, so no one of them is conclusive; we do not claim a demonstrated
failure. Nor do we claim that their agreement in sign constitutes evidence: with
one run per setting, three sub-noise movements in the same measure are equally
consistent with a systematic property of that measure, MME having the widest
dynamic range of our four. We report the pattern because it is the only signal
our design produced, and because it tells a future study where to look, not
because it settles anything. Two things would settle it, and neither is present
here: repeated seeds, which would give the run-to-run variance that our
evaluation-set standard errors do not estimate, and a positive control in which
forgetting is deliberately induced, which would show that these measures move
when a capability is genuinely damaged. Until both exist, our result should be
read as a failure to detect cross-modal forgetting under this protocol, not as a
demonstration of its absence.

\subsection{The flip side: what is learned does not cross modalities}
Localized adaptation is usually reported as good news. We show it has a
symmetric cost. Having acquired Welsh in the text direction, can the model
\emph{act} on Welsh in the visual directions? This is a transfer question, not a
retention one: the base is expected to fail on Welsh because it barely knows it,
so any gain must come from the adapter (protocol in Appendix~\ref{app:xmodal}).

Transfer was mechanistically plausible here, which is why the negative answer is
worth reporting. Janus' image-generation forward pass traverses the same
language-model tensors that the adapter modifies, so a text-side update reaches
the generation pathway by construction; nothing in the architecture prevents a
representation acquired in text from conditioning image tokens. A reader who
accepts that a shared backbone carries shared representations should expect at
least partial transfer. We find instead that the pathway is architecturally open
and functionally unused.

Table~\ref{tab:xmodal} shows the answer is essentially no. Generating images from
Welsh prompts, the adapter improves CLIPScore against the same English reference
only marginally over the base ($10.6\to11.0$, $+0.47\pm0.25$, $20/30$ prompts
improved), which closes just $3\%$ of the English--Welsh gap ($15.8$) and stays
far from English-prompt faithfulness. On Welsh visual question answering there is
no transfer at all: on $28$ balanced yes/no items the base scores $53.6\%$ and
the adapter $50.0\%$, both indistinguishable from chance ($95\%$ Wilson intervals
$[36,70]$ and $[33,67]$), against $100\%$ in English for both. The failure is not
a matter of degree. On the Welsh items the base answers ``yes'' to $11$ of the
$14$ questions whose correct answer is ``no'', while answering $12$ of $14$
``yes'' items correctly: it is not reading the Welsh question at all, and falls
back on a response bias that a balanced item set converts into chance-level
accuracy. The English controls are unchanged, so the probe is sound and the
adapter has damaged nothing. Cross-modal grounding of a new language would
require paired image--text data; text-only post-training cannot supply it. The
adapter is localized to the text stream in exactly the sense that protects the
other modalities, and for exactly the same reason it cannot serve them.

\begin{table}[t]
\caption{\textbf{Cross-modal transfer of the acquired language.} On the English
control columns the adapter matches the base, so the probes are sound. On the
Welsh columns the adapter improves text-to-image faithfulness only marginally and
Welsh VQA not at all; both VQA scores are within the $95\%$ Wilson interval of
chance on $n{=}28$ balanced yes/no items. Text-only acquisition does not reach
the visual pathway.}
\label{tab:xmodal}
\centering\small
\begin{tabular}{lcccc}
\toprule
& \multicolumn{2}{c}{Base} & \multicolumn{2}{c}{Welsh LoRA} \\
\cmidrule(lr){2-3}\cmidrule(lr){4-5}
& English & Welsh & English & Welsh \\
\midrule
(A) text$\to$image, CLIPScore $\uparrow$ ($n{=}30$) & 26.4 & 10.6 & 26.6 & 11.0 \\
(B) VQA accuracy $\uparrow$ ($n{=}28$)              & 100\% & 53.6\% & 100\% & 50.0\% \\
\bottomrule
\end{tabular}
\end{table}

\subsection{How large is the intervention?}
\label{sec:dw}
A sceptical reading of everything above is that nothing is forgotten because
nothing is changed. Table~\ref{tab:dw} rules that out. We summarize each update
$\Delta W=(\alpha/r)BA$ by $\lVert\Delta W\rVert_F/\lVert W\rVert_F$, the size of
the update relative to the weights it modifies, a task-vector magnitude
\citep{ilharco2023task}. It is small but real, $0.6$ to $1.9\%$, and it is
\emph{ordered} in both knobs and both directions: rising with rank, shrinking as
the KL weight grows. These sub-$2\%$ updates already shift behaviour
substantially (poetry perplexity $32.7\to23.8$; visible style transfer,
Figure~\ref{fig:teaser}) while the other modality stays flat. Per-module heatmaps
(Appendix~\ref{app:dw}) show the change is localized: the attention $q/k/v$
projections are nearly untouched and the update concentrates in the MLP, broadly
over \texttt{gate} for text and toward the final layer's \texttt{down} for image.
The global percentage therefore understates the local effect, and the two
directions differ in \emph{where} they write.

\begin{table}[t]
\caption{\textbf{Size of the update} $\lVert\Delta W\rVert_F/\lVert W\rVert_F$
(\%, summed over all backbone projections). It grows smoothly with rank and
shrinks as the KL weight grows, more sharply for text. At matched rank the image
adapter changes the shared backbone more than text.}
\label{tab:dw}
\centering\small
\begin{subtable}{0.52\textwidth}
\centering
\caption{vs.\ rank}
\begin{tabular}{lcccccc}
\toprule
Rank $r$ & 2 & 4 & 8 & 16 & 32 & 64 \\
\midrule
Image (gen.)  & 0.58 & 0.75 & 0.97 & 1.25 & 1.57 & 1.91 \\
Text (poetry) & --- & 0.67 & 0.84 & 1.05 & 1.32 & 1.66 \\
\bottomrule
\end{tabular}
\end{subtable}
\hfill
\begin{subtable}{0.44\textwidth}
\centering
\caption{vs.\ KL weight $\lambda$}
\begin{tabular}{lccccc}
\toprule
$\lambda$ & 0 & 0.25 & 0.5 & 1 & 1.5 \\
\midrule
Image ($r{=}32$) & 1.55 & 1.58 & 1.57 & 1.45 & 1.33 \\
Text ($r{=}16$)  & 1.15 & 1.15 & 1.05 & 0.89 & 0.80 \\
\bottomrule
\end{tabular}
\end{subtable}
\end{table}

\paragraph{An anomaly we report but do not explain.}
General-text perplexity falls consistently \emph{below} the base after text
adaptation, across ranks, KL weights and all three targets (e.g.\
$13.71\to12.07$ for Welsh), and the direction is corroborated by the higher
MAUVE of the adapters' free English. In continual-learning terms this would be
positive backward transfer, but we have no confirmed mechanism. Plausible and
unverified causes include a mild regularization of the shared language head by
the additional text objective, and a reduced tendency to slip into the assistant
persona under free continuation. Distinguishing these requires multi-seed runs
we have not done, so we report it as an observation and build no claim on it.

\section{Discussion}
\paragraph{Localization is a trade, not a free lunch.}
Read together, our two main results are the same phenomenon seen from two sides.
A small, MLP-localized update in the shared backbone is enough to acquire a
poetic register, a medical domain, or a new language, and it leaves the other
capability directions at base level, on the four benchmarks, in free-form text,
and in the image-generation distribution. That same localization is why the
acquired language cannot be used through the visual pathway. Practitioners who
post-train a unified model can take the first half as reassurance and must take
the second half as a design constraint: if a new capability is to be available in
more than one modality, the evidence here says it has to be trained there,
because it will not arrive by itself.

\paragraph{The KL weight controls drift, not reach.}
The KL term behaves as expected, trading acquisition for preservation
monotonically, and it removes the one movement we could induce ($\lambda=0$,
Table~\ref{tab:textkl}). But it is orthogonal to the transfer question: at every
$\lambda$, including $\lambda=0$, Gen-CMMD is flat, so no setting of it makes the
text update reach the image pathway. How far the model drifts and how far what it
learns travels are controlled by different things.

\paragraph{Limitations.}
Our evidence comes from a single unified model, Janus-Pro-7B, whose
understanding and generation encoders are \emph{decoupled} over a shared
backbone; because LoRA acts only on that backbone, the generation pathway is
partly insulated by construction, which may itself be part of why adaptation
stays localized. Our protocol is also a \emph{single} adaptation step from a
pretrained checkpoint, not a sequence of tasks: we measure interference between
capabilities that coexist in one checkpoint, which is the sub-problem of
continual learning that unified multimodal models make acute, but it is not
sequential continual learning, and we use ``transfer'' in that restricted sense
rather than in the task-sequence sense of \citet{lopezpaz2017gem}. Chaining our
three text targets into one adapter and re-measuring each after every step is the
natural extension and the one we would run next. Whether our conclusions hold for
fully-tokenized early-fusion models (e.g.\ Chameleon), where a single token
stream shares all weights, is untested and is the central open question; we could
not include such models here because their generation is either disabled in the
public release (Chameleon) or inference-only with inaccessible weights (Anole),
the very constraint that motivated our choice of Janus. We do not vary the base
model or its scale. All runs are single-seed, so the small movements of
\S\ref{sec:boundary} are suggestive rather than established, and the general-text
improvement is unexplained. Our comparison of anti-forgetting objectives is
scoped to forgetting mitigation; orthogonal adaptation-capacity variants such as
DoRA \citep{liu2024dora} and multi-adapter routing are out of scope. Finally,
free-form quality is assessed distributionally rather than by a human study.

\section{Conclusion}
We measured cross-modal catastrophic forgetting under parameter-efficient
post-training of a unified multimodal model, choosing every adaptation target by
how poorly the base model already handles it so that preservation is certified
against a real intervention. Across two directions, five targets, LoRA ranks $2$
to $128$ and KL weights $\lambda\in[0,1.5]$, the capabilities we did not train
stay at base level within evaluation noise, including under a strong
distributional shift toward a language the base barely modelled, and including in
free-form generation, which likelihood benchmarks do not test. Where a cost
appears at all it is small, confined to grounded visual understanding, and
reachable only at the strongest settings. The same localization that delivers
this stability also confines what is learned: the newly acquired language does
not transfer to the visual pathway, although the architecture leaves that pathway
open. We are explicit that single-seed evidence supports ``we could not detect
forgetting,'' not ``there is none''; the two experiments that would convert one
into the other are repeated seeds and a positive control in which forgetting is
deliberately induced, and both are within reach. Beyond them, the highest-value
extensions are a sequential version of this protocol and replication on a
fully-tokenized early-fusion model, where the insulation our conclusions may
partly rest on is absent.

\bibliographystyle{plainnat}
\bibliography{references}
